\section{Appendix}

\subsection{Source Code Repository} \label{section:app:repo}
The repository is hosted on GitHub and the source code can be found there with the data as well:
\url{https://github.com/Maccen26/hmm-project}

\subsection{Source code \& Envoriment} \label{section:app:env}
This code was written with Python (version 3.11) and \texttt{uv} (version 0.9.7) \cite{astral_uv_docs} was used as a package manager.
The libraries used are listed in Table~\ref{tab:libraries}.

\begin{table}[H]
    \centering
    \begin{tabular}{ll}
        \toprule
        \textbf{Library} & \textbf{Version} \\
        \midrule
        \texttt{dotenv}      & $\geq$ 0.9.9 \\
        \texttt{equinox}     & $\geq$ 0.13.7 \\
        \texttt{jax}         & $\geq$ 0.9.0.1 \\
        \texttt{matplotlib}  & $\geq$ 3.10.8 \\
        \texttt{numpy}       & $\geq$ 1.23.5, $<$ 2.3 \\
        \texttt{optax}       & $\geq$ 0.2.8 \\
        \texttt{pandas}      & $\geq$ 3.0.0 \\
        \texttt{seaborn}     & $\geq$ 0.13.2 \\
        \texttt{statsmodels} & $\geq$ 0.14.6 \\
        \bottomrule
    \end{tabular}
    \caption{Python libraries used in this project.}
    \label{tab:libraries}
\end{table}

The source code is organized in a way such that the directory \texttt{src} contains the source code for the project,
the directory \texttt{data} contains the data used for training and testing the model,
the directory \texttt{tests} contains tests for the source code and
the directory \texttt{drivers} contains scripts for running the model and generating the results used in the report.

\subsection{CO\textsubscript{2} Data} \label{section:app:co2_data}

\begin{figure}[H]
    \centering
    \includegraphics[width=0.85\textwidth]{../results/plots/co2_data.png}
    \caption{Observed CO\textsubscript{2} concentration (ppm) over time.}
    \label{fig:app:co2_data}
\end{figure}

\subsection{Class Diagrams} \label{section:app:diagrams}

\begin{figure}[H]
    \centering
    \includegraphics[width=0.9\textwidth]{../docs/diagrams/01_system_composition}
    \caption{System composition overview. The user-facing \texttt{HMM} class owns an
             \texttt{HMMParams} instance composed of a \texttt{BaseTransition} and a
             \texttt{BaseEmission}. Fitting is delegated to a \texttt{BaseSolver} subclass
             and state inference is handled by \texttt{ForwardAlgorithm}, producing
             \texttt{HMMResults} and \texttt{StateResults}.}
    \label{fig:diagram:composition}
\end{figure}

\begin{figure}[H]
    \centering
    \includegraphics[width=0.85\textwidth]{../diagrams/02_emissions}
    \caption{Emission class hierarchy. \texttt{GaussEmission} models state-dependent
             Gaussian distributions with monotonically ordered means enforced via
             cumulative exponentiation of \texttt{log\_mu\_diff}.
             \texttt{AutoregressiveGaussEmission} extends this with AR($p$) dynamics,
             where coefficients are stored unconstrained (\texttt{phi\_tilde}) and
             mapped to $(-1, 1)$ via $\tanh$ to guarantee stationarity.}
    \label{fig:diagram:emissions}
\end{figure}

\begin{figure}[H]
    \centering
    \includegraphics[width=0.85\textwidth]{../diagrams/03_transitions}
    \caption{Transition class hierarchy. Both implementations store parameters as
             logits and convert them to a row-stochastic matrix via row-wise softmax.
             \texttt{StaticTransition} uses a constant $K \times K$ matrix, while
             \texttt{StaticTransitionHigherOrder} augments the state space to
             $K^{\text{order}}$ states to represent a higher-order Markov chain.}
    \label{fig:diagram:transitions}
\end{figure}

\begin{figure}[H]
    \centering
    \includegraphics[width=0.85\textwidth]{../diagrams/04_inference_output}
    \caption{Inference and output class hierarchy. \texttt{BaseInference.run()} iterates
             \texttt{step()} across the time axis using \texttt{jax.lax.scan}.
             \texttt{ForwardAlgorithm} implements the forward filtering pass and produces
             a \texttt{ForwardOutput} storing the scaling factors $f_t$, filtered state
             probabilities $u_{t|t}$, and predicted state probabilities $u_t$.
             The log-likelihood is computed as $\sum_t \log f_t$.}
    \label{fig:diagram:inference}
\end{figure}

\begin{figure}[H]
    \centering
    \includegraphics[width=0.85\textwidth]{../diagrams/05_solvers}
    \caption{Solver class hierarchy. \texttt{BaseSolver} partitions parameters into
             trainable and frozen subsets and wraps the negative log-likelihood as the
             loss function. The three concrete solvers share the same interface but
             differ in their optimization backend: \texttt{GradientSolver} uses
             first-order gradient descent via \texttt{optax}, \texttt{LBFGSSolver}
             uses quasi-Newton L-BFGS via \texttt{optax}, and \texttt{Minimizer}
             wraps \texttt{jax.scipy.optimize.minimize}.}
    \label{fig:diagram:solvers}
\end{figure}


\subsection{Class Design} \label{section:app:class_design}

The \texttt{HMM} class (\ref{fig:diagram:composition}) is the main entry point of the system.
Its responsibility is to act as a wrapper for the other components and to provide a simple interface for training and inference.
It is dependent on \texttt{Emissions}, \texttt{Transitions}, \texttt{HMMParams}, and \texttt{Solvers}, which are composed together to form a complete HMM.

The \texttt{Emissions} (\ref{fig:diagram:emissions}) and \texttt{Transitions} (\ref{fig:diagram:transitions}) classes are abstract interfaces whose concrete implementations must be provided at runtime.
This design choice is motivated by the structure of HMM variants: the differences between models in this project lie in how the density of an observation given a state is computed (emissions) and how the transition probabilities between states are computed (transitions).
By separating these concerns into abstract interfaces, different HMM models can be composed without changing the rest of the system.

The \texttt{HMMParams} class acts as a container for the model parameters and as a wrapper for the methods implemented in the \texttt{Emissions} and \texttt{Transitions} classes.
This is required by how \texttt{JAX} and \texttt{Equinox} work: parameters must be stored as pytrees to support automatic differentiation.

The \texttt{Inference} class (\ref{fig:diagram:inference}) is an abstract interface responsible for implementing inference algorithms such as the forward algorithm, the backward algorithm, and the Viterbi algorithm.
In this project, only the forward algorithm is implemented, but the design allows easy extension to other inference algorithms without modifying existing code.

The \texttt{Solvers} class (\ref{fig:diagram:solvers}) implements the optimization algorithms and depends on a loss function that takes the \texttt{HMMParams} and \texttt{Inference} classes as arguments.
This dependency injection of the loss function decouples the optimizer from the objective, making it straightforward to change the loss function without altering the optimization algorithm.
